\section*{Capítulo \ref{ch:g12:metamorphic-rocks} --- Las rocas metamórficas}

\begin{solution}{exo:g12:metamorphic-rocks:1}
(a) \omterm{def:g12:metamorphic-rocks:polymorph}{polimorfos}; (b) una \omterm{def:g12:metamorphic-rocks:index-mineral}{isógrada}; (c) una \omterm{def:g12:metamorphic-rocks:contact-regional}{aureola}; (d) una \omterm{def:g12:metamorphic-rocks:foliation}{foliación};
(e) la \omterm{def:g12:metamorphic-rocks:lithostatic}{presión litostática}.
\end{solution}

\begin{solution}{exo:g12:metamorphic-rocks:2}
(1) mármol; (2) cuarcita; (3) pizarra; (4) gneis; (5) corneana.
\end{solution}

\begin{solution}{exo:g12:metamorphic-rocks:3}
$P = \rho g z$: a \qty{10}{km}, $2830 \times 9.81 \times \num{10000} \approx
\qty{2.78e8}{Pa} = \qty{0.278}{GPa}$; a \qty{20}{km}, el doble,
\qty{0.555}{GPa}.
\end{solution}

\begin{solution}{exo:g12:metamorphic-rocks:4}
(1) Cianita (la línea cianita--andalucita está hacia \qty{0.25}{GPa} a
\qty{400}{\celsius}). (2) Andalucita. (3) Sillimanita.
\end{solution}

\begin{solution}{exo:g12:metamorphic-rocks:5}
(1) Pizarra, filita, esquisto, gneis. (2) En el esquisto. (3) El gneis.
\end{solution}

\begin{solution}{exo:g12:metamorphic-rocks:6}
(1) $2830 \times 9.81 \times \num{35000} \approx \qty{9.72e8}{Pa}$, es decir,
$\num{9.72e8}/101325 \approx \num{9600}$ atmósferas. (2) $1025 \times 9.81
\times 4000 \approx \qty{4.02e7}{Pa} = \qty{40.2}{MPa}$. (3) $972/40.2
\approx 24$ veces.
\end{solution}

\begin{solution}{exo:g12:metamorphic-rocks:7}
(1) $z = \num{0.375e9}/(2830 \times 9.81) \approx \qty{13.5}{km}$.
(2) $(504 - 10)/13.5 \approx \qty{36.6}{\celsius}$ por kilómetro. (3) No:
la roca tendría que encontrarse exactamente en un punto del diagrama; esas
rocas son raras.
\end{solution}

\begin{solution}{exo:g12:metamorphic-rocks:8}
(1) $23/3 \approx \qty{8}{mm}$. (2) Unos 8 millones de años. (3) La roca
permaneció en profundidad, a alta temperatura, durante millones de años.
\end{solution}

\begin{solution}{exo:g12:metamorphic-rocks:9}
(1) Cuatro: las \omterm{def:g12:metamorphic-rocks:index-mineral}{isógradas} del granate, la estaurolita, la cianita y la
sillimanita, en ese orden. (2) Hacia el oeste. (3) En el extremo occidental,
en la zona de la sillimanita. Sus rocas cruzaron la línea
cianita--sillimanita: enterradas a unos \qty{20}{km} (\qty{0.56}{GPa}),
estaban a más de unos \qty{590}{\celsius}.
\end{solution}

\begin{solution}{exo:g12:metamorphic-rocks:10}
(1) El fango se convierte en corneana, la caliza en mármol y la arenisca en
cuarcita. (2) Se calentaron sin ser comprimidas: ninguna presión dirigida
alineó sus minerales. (3) Las rocas que lo rodean ya están calientes, así que
el magma apenas las calienta más.
\end{solution}

\begin{solution}{exo:g12:metamorphic-rocks:11}
(1) Perpendiculares al empuje: sus caras planas miran al este y al oeste, así
que las escamas quedan en planos norte--sur. (2) Vertical, de dirección
norte--sur. (3) Los granos de calcita y de cuarzo no son ni planos ni
alargados, así que, aunque recristalicen bajo compresión, no hay nada que
alinear.
\end{solution}

\begin{solution}{exo:g12:metamorphic-rocks:12}
(1) $z = (T - 10)/30$, así que $P = 0.0278\,(T - 10)/30 \approx
0.000927\,(T - 10)$. (2) $0.000927\,(T - 10) = 0.577 + 0.00209\,(T - 600)$
da $0.001163\,T \approx 0.668$, así que $T \approx \qty{574}{\celsius}$, $P
\approx 0.000927 \times 564 \approx \qty{0.523}{GPa}$ y $z = (574 - 10)/30
\approx \qty{18.8}{km}$. (3) Cianita por encima de esa profundidad,
sillimanita por debajo.
\end{solution}

\begin{solution}{exo:g12:metamorphic-rocks:13}
(1) $10 + 10 \times 35 = \qty{360}{\celsius}$. (2) $10 + 30 \times 35 =
\qty{1060}{\celsius}$. (3) La presión a \qty{35}{km} es de unos
\qty{0.97}{GPa}; a \qty{360}{\celsius}, ese punto está en pleno campo de la
cianita. El segundo punto queda muy a la derecha del diagrama, en el campo de
la sillimanita, a temperaturas a las que las rocas de la corteza empiezan a
fundirse. (4) Necesitan una presión alta a baja temperatura, una combinación
que solo proporciona la corteza fría arrastrada hacia abajo en una \omterm{def:g11:subduction-zones:subduction}{zona de
subducción}.
\end{solution}

\begin{solution}{exo:g12:metamorphic-rocks:14}
(1) Cianita, $1000/3.61 \approx \qty{277}{cm^3}$; sillimanita, $1000/3.24
\approx \qty{309}{cm^3}$; andalucita, $1000/3.15 \approx \qty{317}{cm^3}$.
(2) $(317 - 277)/317 \approx 0.127$, alrededor del \qty{13}{\percent}. (3) A
alta presión se favorece la disposición más compacta de los átomos; la
cianita empaqueta los mismos átomos en el menor espacio.
\end{solution}

\begin{solution}{exo:g12:metamorphic-rocks:15}
(1) \qty{0.555}{GPa} a \qty{20}{km}, \qty{0.694}{GPa} a \qty{25}{km}.
(2) $T = 600 + (P - 0.577)/0.00209$: unos \qty{590}{\celsius} a
\qty{0.555}{GPa} y \qty{656}{\celsius} a \qty{0.694}{GPa}. A
\qty{700}{\celsius}, las rocas estaban más calientes que ambos valores, en el
campo de la sillimanita: es compatible. (3) $(700 - 10)/20 \approx \qty{34.5}{\celsius}$ y
$(700 - 10)/25 \approx \qty{27.6}{\celsius}$ por kilómetro, cerca del
gradiente típico. (4) De $\qty{20}{km}$ a $\qty{25}{km}$ en 40 millones de años:
de 0.50 a \qty{0.63}{km} por millón de años, es decir, de 0.50 a
\qty{0.63}{mm} al año.
\end{solution}

\begin{solution}{pb:g12:metamorphic-rocks:1}
\textbf{1.} Está foliado, y sus micas son lo bastante grandes para verse. Las
micas, moscovita y biotita, forman la \omterm{def:g12:metamorphic-rocks:foliation}{esquistosidad}.

\textbf{2.} \ce{Al2SiO5}; son \omterm{def:g12:metamorphic-rocks:polymorph}{polimorfos}.

\textbf{3.} Los volúmenes de un kilogramo están en la razón $3.15/3.61 \approx
0.873$: el de la cianita es alrededor de un \qty{13}{\percent} menor. La
presión favorece la cianita.

\textbf{4.} La zona de la cianita. En su descenso cruzó las \omterm{def:g12:metamorphic-rocks:index-mineral}{isógradas} de la
biotita, el granate, la estaurolita y la cianita.

\textbf{5.} Unos 8 millones de años.

\textbf{6.} $P = 0.577 + 0.00209 \times 0 = \qty{0.577}{GPa}$.

\textbf{7.} \qty{0.70}{GPa} está por encima de \qty{0.577}{GPa}, del lado de
la cianita de la línea cianita--sillimanita; también está muy por encima del
campo de la andalucita, que a \qty{600}{\celsius} queda por debajo de unos
\qty{0.23}{GPa}.

\textbf{8.} \qty{0.577}{GPa}.

\textbf{9.} $0.70 = 0.577 + 0.00209\,(T - 600)$, así que $T = 600 +
0.123/0.00209 \approx \qty{659}{\celsius}$.

\textbf{10.} \qty{7.0}{kbar}; $\num{7.0e8}/101325 \approx \num{6900}$
atmósferas.

\textbf{11.} $\rho g = 2830 \times 9.81 \approx \qty{27800}{Pa}$ por metro,
\qty{27.8}{MPa} por kilómetro.

\textbf{12.} $z = \num{0.70e9}/\num{27762} \approx \qty{25200}{m}$, unos
\qty{25.2}{km}.

\textbf{13.} $\num{0.577e9}/\num{27762} \approx \qty{20.8}{km}$.

\textbf{14.} $\num{0.70e9}/(2700 \times 9.81) \approx \qty{26.4}{km}$, alrededor
de un \qty{5}{\percent} más profundo ($26.4/25.2 \approx 1.048$).

\textbf{15.} En la corteza: \qty{25}{km} es menos incluso que el espesor de
una corteza normal.

\textbf{16.} $(600 - 10)/25.2 \approx \qty{23.4}{\celsius}$ por kilómetro.

\textbf{17.} Es mucho más caliente que una \omterm{def:g11:subduction-zones:subduction}{zona de subducción} (menos de
\qty{10}{\celsius} por kilómetro) y algo más fría que una corteza típica
(\qty{30}{\celsius}), lejos de una zona volcánica (de 40 a \qty{50}{\celsius}):
una roca enterrada a gran profundidad en la corteza engrosada de un cinturón
montañoso, por \omterm{def:g12:metamorphic-rocks:contact-regional}{metamorfismo regional}.

\textbf{18.} $10 + 30 \times 25.2 \approx \qty{766}{\celsius}$. La línea
cianita--sillimanita, prolongada, está a $0.577 + 0.00209 \times 166
\approx \qty{0.92}{GPa}$ a esa temperatura; \qty{0.70}{GPa} queda por debajo,
en el campo de la sillimanita. Una \omterm{def:g11:lithosphere-asthenosphere:geotherm}{geoterma} típica habría formado sillimanita.

\textbf{19.} $\qty{25.2}{km} = \qty{2.52e7}{mm}$; a \qty{0.5}{mm} al año,
unos 50 millones de años.

\textbf{20.} Al ascender, la roca se enfrió, así que sus reacciones se
frenaron, y el agua que necesitan las reacciones inversas había sido
expulsada durante el descenso.

\textbf{21.} El esquisto con cianita cristalizó a unos \textbf{\qty{25}{km}
de profundidad}, a unos \qty{600}{\celsius}.
\end{solution}
